本章将介绍各种不同的架构方法或架构风格。每一节都会讨论一种不同的软件设计方法，介绍其优缺点，并说明何时以及如何应用它才能获得收益。本章首先会比较有状态与无状态的架构。之后，我们将介绍单体系统、面向服务的设计和微服务。接着，我们会从不同角度考察各种架构风格，依次介绍基于事件的架构、事件溯源、分层系统、模块化设计，最后是六边形架构。

完成本章后，你将熟悉以下主题：

\begin{itemize}
\item 在有状态与无状态方法之间做出选择
\item 理解单体（应用）
\item 理解服务与微服务
\item 探索基于事件的架构
\item 探索事件溯源模式
\item 理解分层架构
\item 学习基于模块的架构
\item 介绍六边形架构
\end{itemize}

\section{技术要求}

你需要了解什么是软件服务。本章的代码可以在以下 GitHub 页面找到：\href{https://github.com/PacktPublishing/Software-Architecture-with-Cpp-2E/tree/main/Chapter02}{https://github.com/PacktPublishing/Software-Architecture-with-Cpp-2E/tree/main/Chapter02}。

\section{在有状态与无状态方法之间做出选择}

\label{idx:073}
有状态和无状态是实现 API 的两种截然相反的方式。两者各有优缺点。

\label{idx:074}
顾名思义，有状态软件的行为取决于其内部状态。以\emph{图 2.1}所示的 Web 服务为例。如果服务能记住自身状态，它的使用方就可以在每次请求中发送更少的数据，因为服务会记住这些请求的上下文。然而，在请求大小和带宽上的节省，会在 Web 服务这一侧带来隐性成本。如果用户同时发送大量请求，服务就必须同步自己的工作。由于多个并发请求都可能改变状态，如果缺乏同步，就可能导致数据竞争，从而引发未定义或不可预测的行为。

\myGraphic{0.920}{images/B18456_02_01.png}{图 2.1：有状态与无状态的 API}

\begin{myTip}{提示}
想查看此图像的高分辨率版本？请在下一代 Packt Reader 中打开本书，或在 PDF/ePub 版本中查看。

\textbf{下一代 Packt Reader} 以及本书的 \textbf{免费 PDF/ePub 副本} 随购买一同提供。扫描二维码或访问 \href{https://packtpub.com/unlock}{https://packtpub.com/unlock}，然后使用搜索栏按书名查找本书。请仔细核对显示的版本，确保拿到的是正确的那一版。

\begin{center}\includegraphics[width=0.250\linewidth]{images/01.png}\end{center}
\end{myTip}

\label{idx:075}
\label{idx:076}
\label{idx:077}
\label{idx:078}
但若服务是无状态的，那么发往它的每个请求都必须包含成功处理该请求所需的全部数据。这意味着请求会变得更大、占用更多带宽，但另一方面，它能让服务获得更好的性能和扩展性。如果你熟悉函数式编程，可能会觉得无状态服务很直观。处理每个请求可以理解为调用一个纯函数，在给定相同输入时返回相同输出。事实上，无状态编程带来的许多优势都源自其函数式编程的根。可变状态是并发代码的大敌。函数式编程依赖不可变的值，即便这意味着要创建副本，而不是修改现有对象。正因如此，每个线程都可以独立工作，不可能发生数据竞争。我们将在\emph{第 5 章}\emph{《利用 C++ 语言特性》}中更详细地讨论函数式编程。

由于不存在竞态条件，也就不需要锁来保护共享状态、避免它被多个线程同时访问。这能带来巨大的性能提升，因为维护和访问会话数据不再有任何开销。不用锁还意味着你不再需要处理死锁——当两个线程共享同一资源、互相阻止对方访问该资源而停止运行时，就会发生死锁。使用纯函数还意味着代码更容易调试，因为没有任何副作用，也就是程序或环境状态的改变没有反映在函数输出中。反过来说，没有副作用对编译器也有好处，因为优化无副作用的代码要容易得多，也就能更激进地进行优化。这并不意味着你应该总是采用无状态，因为选择取决于许多因素，包括项目需求。如果你需要支持会话并跟踪用户状态，有状态的方法可能更合适。如果性能和可扩展性更重要，那么无状态的方法会更好。

以类为例。假设你要建模一个 \texttt{Consultant} 类，那么该类包含顾问的姓名、联系方式、时薪、当前和以往的项目等字段是合理的。它有状态是自然而然的事。现在，假设你需要计算他们工作所得的报酬。你应该创建一个 \texttt{Payment\allowbreak{}Calculator} 类吗？还是应该添加一个成员函数或自由函数来计算它？

如果采用类的方案，你应该把 \texttt{Consultant} 实例作为构造函数参数还是方法参数传入？这个类是否应该拥有津贴之类的属性？

\label{idx:079}
\label{idx:080}
\label{idx:081}
添加一个计算报酬的成员函数会破坏\textbf{单一职责原则}（\textbf{SRP}），因为这样一来，该类就承担了两项职责：计算报酬和存储顾问的数据（状态）。这意味着，相比拥有这类混合职责的类，更应该为此引入一个自由函数或单独的类。

定义 SRP 的另一种方式如下：

\begin{myQuote}
将因相同原因而变化的事物聚集在一起。将因不同原因而变化的事物分离开来。

—— Robert C. Martin，《单一职责原则》，The Clean Coder Blog，2014

\end{myQuote}

也就是说，如果我们修改某个操作的逻辑，就应该只在一个地方修改。换句话说，SRP 不仅要求 \emph{能拆就拆}，还要求不要拆得过头——\emph{不要把相关的事物拆开}。不要无谓地把事情复杂化。

这样的类究竟该不该包含状态？下面来讨论 \texttt{Payment\allowbreak{}Calculator} 类的几种不同做法。

一种做法是公开计算所需的属性：

\begin{codebox}[cpp]
class PaymentCalculator final {
public:
    double calculate() const { return hours_ * netHourlyRate_ * taxPercentage_; }
    void setHours(double hours) { hours_ = hours; }
    void setHourlyRate(double rate) { netHourlyRate_ = rate; }
    void setTaxPercentage(double tax) { taxPercentage_ = tax; }
private:
    double hours_{};
    double netHourlyRate_{};
    double taxPercentage_{};
};
\end{codebox}

\label{idx:082}
\label{idx:083}
这种做法有两个缺点。第一，它不是线程安全的；在没有锁的情况下，\texttt{Payment\allowbreak{}Calculator} 类的单个实例无法在多个线程中使用。第二，一旦计算变得更为复杂，该类很可能会开始重复 \texttt{Consultant} 类中的更多字段。

为了消除这种重复，我们可以改造这个类，让它存储一个 \texttt{Consultant} 实例，像这样：

\begin{codebox}[cpp]
class PaymentCalculator final {
public:
    double calculate() const { return {}; };
    void setConsultant(const Consultant &c) { consultant_ = c; };
    void setTaxPercentage(double tax) { taxPercentage_ = tax; };
private:
    std::reference_wrapper<const Consultant> consultant_;
    double taxPercentage_{};
};
\end{codebox}

这种做法仍有不是线程安全的缺点。我们还能做得更好吗？事实证明，我们可以通过让该类变为无状态来使其线程安全：

\begin{codebox}[cpp]
class PaymentCalculator {
public:
    static double calculate(const Consultant &c, double taxPercentage) { return {}; };
};
\end{codebox}

\label{idx:084}
\label{idx:085}
如果没有状态需要管理，那么你是决定创建自由函数（也许放在一个单独的命名空间中），还是像前一段代码那样把它们归为类的静态函数，其实都无所谓。就类而言，区分数值（实体）类型和操作类型是很有用的，因为把二者混在一起会导致违反 SRP。

\subsection{无状态与有状态的服务}

我们讨论过的这些适用于类的原则，同样可以映射到更高层的概念上，例如微服务。

\label{idx:086}
有状态的服务是什么样子？以\textbf{文件传输协议}（\textbf{FTP}）为例。如果它不是匿名的，就需要用户提供用户名和密码来创建会话。服务器会存储这些数据，以确认用户仍然处于连接状态，因此它在不断地存储状态。用户每次更改工作目录，状态都会更新。用户所做的每一次更改都会反映为状态的改变，即使断开连接时也是如此。拥有有状态的服务意味着，取决于状态的不同，两个看起来完全相同的 \texttt{GET} 请求可能返回不同的结果。如果服务器丢失了状态，你的请求甚至可能无法被正确处理。

有状态的服务还可能面临会话不完整或事务未完成的问题，并带来额外的复杂性。会话应该保持打开多久？如何验证客户端是否已崩溃或断开连接？何时应该回滚已做的更改？虽然你可以为这些问题找到答案，但通常更简单的做法是让服务的使用方以动态、\emph{智能}的方式与服务通信。既然它们会自行维护某种状态，那么让服务也维护状态不仅没有必要，往往还是一种浪费。

\label{idx:087}
无状态服务则采取相反的做法。每个请求都必须包含成功处理它所需的全部数据，因此两个相同的幂等请求（例如 \texttt{GET}）会产生相同的响应。这里假设服务器上存储的数据没有变化，但数据未必等同于状态。真正重要的只是每个请求都是自包含的。

\label{idx:088}
无状态是现代互联网服务的根本特性。\textbf{超文本传输协议}（\textbf{HTTP}）就是无状态的，许多服务 API（例如 X 的 API）也是无状态的。X 的 API 所依赖的\textbf{表述性状态转移}（\textbf{REST}）在设计上就是功能无状态的。REST 背后的整个理念是：处理请求所需的所有状态都必须在请求内部传递。如果不是这样，那就不能说你拥有一个 RESTful 服务。不过，出于实际需要，这条规则也有一些例外。

如果你在构建一个在线商店，可能希望存储与客户相关的信息，例如他们的订单历史和收货地址。客户一侧的客户端可能会存储认证 cookie 或令牌，而服务器则可能会把一些用户数据存储在数据库中。

在基于会话的认证中，cookie 存储在客户端，但指向在服务器上管理的会话。这意味着双方都需要参与。而在基于令牌的认证中，服务器创建一个令牌，其中包含识别用户所需的全部信息，并将其发送给客户端。此后，客户端会在每次请求中带上这个令牌。服务器只需验证令牌的有效性，而无须在中间跟踪任何会话数据，这让事情变得更简单。

把会话保留在服务器端，对服务而言可能是一种受限的做法，原因有几点：它们增加了许多本可避免的复杂性，让缺陷更难复现，而且最重要的是，它们无法扩展。如果你想将负载分布到另一台服务器上，很可能会在随负载复制会话以及在各服务器之间同步会话时遇到麻烦。基于会话的认证最适合这样的场景：你需要实时会话控制或集中式管理，或者可扩展性不是首要考虑因素。

\label{idx:089}
与 cookie 方案不同，令牌方案在服务器端没有会话状态，服务器只负责创建和验证用于标识用户的令牌，从而实现更具扩展性的解决方案。\emph{图 2.2} 展示了服务器端基于会话的认证与基于令牌的认证之间的区别：

\myGraphic{0.920}{images/B18456_02_02.png}{图 2.2：服务器端基于会话的认证与基于令牌的认证}

\label{idx:090}
这意味着，如果你希望采用有状态的架构，就需要有充分的理由。以 FTP 协议为例。它必须在客户端和服务器端都复制这些更改。用户只向一个特定的服务器认证，以完成单次状态的数据传输。把它与 Dropbox 之类的服务相比较——在那里，数据常常在用户之间共享，文件访问通过 API 被抽象掉——就能明白为什么无状态模型更适合这种情况。

\section{理解单体（应用）}

开发应用时最简单的架构风格是单体式。正因如此，许多项目一开始就采用这种风格。单体应用就是一个大块整体，这意味着应用中功能上可区分的部分——例如处理 I/O、数据处理和用户界面——全都交织在一起，而不是分置于不同的架构组件中。单体应用具有以下优点：

\begin{itemize}
\label{idx:091}
\item 部署单体应用可能比部署多组件应用更容易，因为只需要部署一个东西
\item 它们也可能更容易测试，因为端到端测试只需启动一个组件
\item 系统不同部分之间的集成通常更容易
\item 只需在负载均衡器后面增加更多实例，就能扩展解决方案
\end{itemize}

既然有这么多优点，为什么你不采用这种架构风格呢？事实证明，它也有许多缺点：

\begin{itemize}
\label{idx:092}
\item 启动单体应用可能比启动更小、更集成的服务要慢得多。请注意，启动时间取决于所用的技术和系统的复杂度。此外，无论你在应用中改动了什么，你可能都不喜欢它迫使你重跑所有测试并一次性重新部署整个应用。（我们会在\emph{第 10 章}\emph{《编写可测试的代码》}中更详细地探讨测试。）现在，假设你的一位开发人员在应用中引入了资源泄漏。如果这段泄漏的代码被反复执行，就不只是拖垮应用的一部分，而是整个应用。
\label{idx:093}
\item 更重要的是，你开发单体应用的时间越长，维护它时遇到的问题就越多。让这类应用的内部保持松散耦合是一项挑战，因为在它的各个模块之间再添加一个依赖实在太容易了。随着这类应用不断增长，它会越来越难以理解，因此由于复杂性增加，开发过程很可能随时间推移而变慢。在开发单体时，维护设计驱动开发（design-driven development）的限界上下文也可能很困难。但许多公司，例如 Instagram，都设法驾驭住了这类应用。Instagram 拥有一个数百万行的单体 Python 应用，并制定了严格的开发文化规范。
\item 所提供的可扩展性在理论上听起来不错，但如果你的应用包含资源需求各不相同的模块呢？如果你只需要扩展应用中的某一个模块呢？缺乏模块化——这是单体系统固有的特性——正是与这种架构相关的许多缺陷的根源。
\item 如果你热衷于在项目中使用最前沿的技术，那么单体风格并不适合你。因为你现在需要一次性迁移整个应用，更新库或框架可能会变得很复杂。
\end{itemize}

上述说明让人觉得单体架构只适用于简单小巧的应用，但事实并非完全如此。在某些情况下，采用它确实可能是个好主意。如果你关注性能，那么与微服务相比，采用单体有时能帮助你在延迟或吞吐量方面从应用中榨取更多性能。

进程间通信总会带来一些开销，而单体应用则无须承担这部分开销。如果你对相关测量感兴趣，可以参阅 Robin Flygare 和 Anthon Holmqvist 的论文 \emph{《单体系统与基于微服务的系统之间的性能特征》}。

\section{理解服务与微服务}

由于单体架构的种种缺点，其他方法应运而生。一种常见的思路是把你的解决方案拆分为多个相互通信的服务。这样你就可以把开发工作分配给不同的团队，每个团队负责一个单独的服务。与单体架构风格不同，每个团队工作的边界都很清晰。

\subsection{面向服务的架构}

\label{idx:094}
\textbf{面向服务的架构}，简称 \textbf{SOA}，意味着将业务功能模块化，并作为单独的服务提供给消费方应用使用。每个服务都应有一个自描述的接口，并隐藏一切实现细节，例如内部架构、所使用的技术或编程语言。这让多个团队可以按自己喜欢的方式开发服务，也就是说，在底层，每个团队都可以采用最适合自身需求的方案。如果你有两个开发团队，一个擅长 C\#，另一个擅长 C++，他们可以开发出两个能够轻松相互通信的服务。

\label{idx:095}
SOA 的倡导者提出了一份宣言，将以下内容置于优先地位：

\begin{itemize}
\item 业务价值高于技术策略
\item 战略目标高于特定项目的收益
\item 内在的互操作性高于定制的集成
\item 共享服务高于专用实现
\item 灵活性高于优化
\item 渐进式改进高于追求初始的完美
\end{itemize}

\label{idx:096}
\label{idx:097}
尽管这份宣言并未把你绑定到特定的技术栈、实现方式或服务类型上，但最常见的两类服务是\textbf{简单对象访问协议}（\textbf{SOAP}）和 REST。除此之外，近来还有第三类越来越流行：\textbf{gRPC 远程过程调用}（\textbf{gRPC}）。

在为项目选择 SOA 之前，你需要考虑实现的复杂度和系统的测试难度。此外，采用这种架构可能会导致网络负载增加，并使某些服务出现延迟。

\subsection{微服务}

\label{idx:098}
顾名思义，微服务是一种架构模式和组织模式，它把应用拆分为一组松散耦合的服务，这些服务通过网络通信协议相互通信。微服务模式与 Unix 哲学相似，后者主张一个程序只应有一个用途。按照 Unix 哲学，复杂的问题要通过把这类程序组合成 Unix 管道来解决。与此类似，基于微服务的系统由许多微服务和支撑服务组成。

为了达到期望的质量属性，微服务依赖一些成熟的模式，例如隔离故障组件以确保容错；但为了获得高效的架构，你必须确定对 API 网关、服务注册表、负载均衡、容错、监控、配置管理，当然还有所要使用的技术栈等要素的处理方式。

\emph{图 2.3} 展示了单体、面向服务和微服务架构的示意图，以帮助区分它们：

\myGraphic{0.920}{images/B18456_02_03.png}{图 2.3：单体、面向服务与微服务架构的对比}

\label{idx:099}
下面先来概述这种架构风格的特征。

\subsubsection{微服务的特征}

\label{idx:100}
由于微服务风格出现得比较晚，目前还没有统一的定义。Martin Fowler 在其题为 \emph{Microservices} 的文章中概括了微服务的若干核心特征，我们将在下面逐一介绍：

\begin{itemize}
\item 每个服务都应当是可独立替换、可独立升级的组件。这带来更简单的部署以及服务之间的松耦合，而不是像单体应用中的组件那样只是一个个库。在后一种情况下，替换其中一个库往往需要重新部署整个应用。
\item 每个服务都应当由一个跨职能团队开发，并聚焦于某项具体的业务能力。不妨想想 \textbf{康威定律}：
\end{itemize}

\begin{myQuote}
任何设计系统（广义上的系统）的组织，其产出的设计结构都会是该组织沟通结构的翻版。

——Melvyn Conway，How Do Committees Invent?，1968 年

\end{myQuote}

\label{idx:101}
如果没有跨职能团队，最终就会形成软件孤岛。随之而来的沟通缺失意味着，你必须不断跨越重重障碍才能成功交付。

\begin{itemize}
\item 每个服务都应当是一个产品，在其整个生命周期内由开发团队拥有。这与项目式思维形成对比：在项目式思维中，你开发完软件就把它交出去，让别人去维护。
\label{idx:102}
\item 服务应当拥有智能端点并使用哑管道，而不是反过来。这与传统服务形成对比：传统服务往往依赖 \textbf{企业服务总线}（\textbf{ESB}）的逻辑，由 ESB 管理消息路由并按业务规则转换消息。在微服务中，你把逻辑存放在服务内部来实现内聚，并避免与消息传递组件耦合。使用消息队列（例如 ZeroMQ）有助于实现这一目标。
\item 服务应当以去中心化的方式治理。单体通常采用某一种特定的技术栈编写。当把它拆分成微服务时，每个微服务都可以选择最适合自身具体需求的技术。保证每个微服务 7×24 小时运行并对其进行治理，由负责该特定服务的团队来完成，而不是由一个中心化的部门来做。Amazon 和 Netflix 等公司就采用这种做法，它们发现，让开发者对自己服务在生产环境中的无故障运行负责，有助于确保高质量。
\item 服务应当以去中心化的方式管理自己的数据。与其让所有服务共用同一个数据库，每个微服务都可以选择最适合自身需求的数据库。去中心化的数据可能在处理更新时带来一些挑战，但它能带来更好的扩展性。正因如此，微服务往往以无事务的方式协同，并提供最终一致性：只要分布式系统中的所有节点在一定时间后达到相同状态，就允许出现暂时的不一致。在分布式系统中，数据可能存放在不同的服务器上，往往还位于不同的数据中心。因此，这些数据可能无法在系统的所有节点之间即时同步。
\label{idx:103}
\item 服务所使用的基础设施应当自动管理。为了高效地应对数十个微服务，你需要建立持续集成与持续交付，否则部署服务将是一场噩梦。让所有测试自动化运行，能帮你省下大量时间和麻烦。在此基础上再实施持续部署，就能缩短反馈回路，也能让你的客户更快用上新功能。
\item 微服务应当对它所依赖的其他服务发生故障有所准备。在一个有如此多活动部件的分布式部署环境中，其中一些时不时出故障是很正常的。你的服务应当能够优雅地处理这类故障。断路器（Circuit Breaker）和舱壁（Bulkhead）等模式（本书后文会介绍）有助于做到这一点。要让架构具备韧性，还必须能够高效地把故障服务重新拉起，甚至提前预知它们将要崩溃。为此，对延迟、吞吐量和资源使用情况进行实时监控至关重要。去了解一下混沌工程工具，例如 Chaos Mesh、Chaos Monkey、Gremlin、Harness Chaos Engineering 和 LitmusChaos，它们对于构建韧性架构极为宝贵。
\item 基于微服务的架构应当随时准备持续演进。你在设计微服务以及它们之间的协作方式时，应当使其能够轻松替换单个微服务，有时甚至替换整组微服务。
\end{itemize}

要恰当地设计这些服务并不容易，尤其因为过去存在于某个较大模块代码中的一部分复杂性，如今可能以服务之间复杂的通信方案的形式出现。这意味着，相比传统服务或单体方案，架构师的经验与技能组合发挥着更为重要的作用。

\label{idx:104}
除此之外，许多（并非全部）微服务还具备以下一些其他特征：

\begin{itemize}
\item 使用彼此独立的进程，通常通过网络协议进行通信
\item 使用与技术无关的协议（例如 HTTP、AMQP、gRPC 或 MQTT）
\item 让服务保持小型化、细粒度且相互独立，并让进程间通信的开销保持在较低水平
\end{itemize}

现在，你应当已经很好地理解了基于微服务的系统具备哪些特征，下面来看看这种方法与其他架构风格相比如何。

\subsubsection{微服务的优势与劣势}

\label{idx:105}
微服务架构中服务的小型化意味着它们开发、部署和理解起来都更快。由于各服务是彼此独立构建的，编译它们新版本所需的时间可以大幅缩短。得益于此，采用这种架构风格时更容易进行快速原型设计和开发。这反过来又能缩短交付周期，也就是说，业务需求可以被更快地引入和评估。

基于微服务的方法还有一些其他优势，如下所示：

\begin{itemize}
\item \textbf{模块化}，这是该架构风格与生俱来的特性
\item \textbf{更好的可测试性}
\item 替换系统组成部分（例如单个服务、数据库、消息代理或云提供商）时的\textbf{灵活性}
\item \textbf{与遗留系统的集成}：无需迁移整个应用，只需迁移当前需要开发的部分
\item \textbf{支持分布式开发}：相互独立的开发团队可以并行开发多个微服务
\item \textbf{可扩展性}：一个微服务可以独立于其他微服务进行扩展
\end{itemize}

\label{idx:106}
另一方面，微服务也有一些劣势：

\begin{itemize}
\item 它们需要成熟的 DevOps 方法，并依赖 CI/CD 自动化
\item 它们更难调试，需要更好的监控和分布式追踪
\item 对于较小的应用，进程间通信带来的额外开销（就辅助服务而言）可能超过其带来的收益
\end{itemize}

下面来看看这种架构广受青睐的主要原因之一——可扩展性。

\subsubsection{扩展微服务}

\label{idx:107}
微服务的扩展方式与单体应用不同。在单体中，全部功能由单个进程处理，扩展通常通过横向或纵向扩展来实现。横向扩展（复制）是指在不同的机器上复制应用，以提升系统的容量；纵向扩展则是为了同样的目的而增加处理器、内存、磁盘等资源。扩展单体不需要考虑哪些功能被大量使用、哪些功能不需要额外的资源。

扩展立方体（scale cube）模型对于理解微服务的设计原则很有用。如 \emph{图 2.4} 所示，该模型涵盖三种相互正交的提升应用性能的方式：横向复制（镜像）、功能分解（微服务）和数据分区（分片）。

\myGraphic{0.746}{images/B18456_02_04.png}{图 2.4：扩展立方体}

\label{idx:108}
横向复制是指克隆服务并在它们之间分发流量，以应对流量高峰。这样一来，系统就能在预期的时限内响应请求。

功能分解是指把单体应用拆分成更小、更专注、松耦合、高内聚、可独立部署且更易于管理的微服务。借助微服务，可以只复制那些负载繁重的服务，而不会影响整个应用。

数据分区通常用于扩展数据库，以提升其可用性和性能，因为每个实例负责一个数据子集。分区是把数据库内的数据拆分成若干子集，而分片则是把数据分布到多个数据库或服务器上。不要把分片与复制混淆：在复制中，数据库实例并不是共享存储的各个部分，而是彼此的副本。

\subsubsection{向微服务迁移}

\label{idx:109}
大多数公司都有某种现存的单体代码，它们不想立刻用微服务重写，但又希望向这种架构过渡。在这种情况下，可以逐步引入微服务：不断增加与单体交互的服务。你可以把新功能做成微服务，也可以直接从单体中切出某些部分，把它们改造成微服务。

有关迁移到微服务的更多细节可参见 \emph{第 14 章}《\emph{分布式系统架构}》。

\section{探索基于事件的架构}

\label{idx:110}
基于事件的系统围绕事件处理展开，事件被定义为系统状态的变化。其中有的组件负责生成事件，有的通道用于传播事件，还有的监听者对事件作出反应，并可能触发新的事件。这种风格提倡异步和松耦合，从而提升性能、可扩展性和部署的便捷性。

有这些优势，也存在一些挑战。其中之一是构建此类系统的复杂性。必须让所有队列具备容错能力，以免事件在处理过程中丢失。以分布式方式处理事务本身也是一项挑战。使用关联 ID（correlation ID，附加到请求或消息上的唯一标识符）来跟踪进程之间的事件，再配合监控技术，能为你省下数小时的调试与冥思苦想。

基于事件的系统的例子包括流处理器和数据集成，以及那些追求低延迟或高可扩展性的系统。

下面来讨论此类系统中常用的拓扑结构。事件驱动架构的两种主要拓扑是基于代理（broker）的和基于中介者（mediator）的。这两种拓扑的区别在于事件在系统中的流动方式。

当处理一个需要执行多个可独立完成的任务或步骤的事件时，最适合使用中介者拓扑。最初产生的所有事件都会进入中介者的事件队列。中介者知道处理该事件需要做哪些事，但它并不亲自执行这些逻辑，而是通过各个处理器的事件通道把事件分发给相应的事件处理器。

\label{idx:111}
\label{idx:112}
\label{idx:113}
如果这让你联想到业务流程的流转方式，那说明你的直觉不错。你可以通过\textbf{业务流程管理}（\textbf{BPM}）或\textbf{业务流程执行语言}（\textbf{BPEL}）来实现这种拓扑。不过，你也可以使用 Apache Camel 或 Mule ESB 之类的技术来实现它。

\myGraphic{0.920}{images/B18456_02_05.png}{图 2.5：中介者拓扑}

而代理则是一个轻量级组件，它包含所有队列，但不会编排事件的处理过程。它可以要求接收方订阅特定类型的事件，然后只需把接收方感兴趣的事件转发给它们。许多消息队列都依赖代理，例如 ZeroMQ，它是用 C++ 编写的，目标是零浪费和低延迟。

\myGraphic{0.920}{images/B18456_02_06.png}{图 2.6：代理拓扑}

\label{idx:114}
现在你已经了解了基于事件的系统中两种常用的拓扑，接下来我们来学习一种以事件为核心的强大架构模式。

\section{探索事件溯源模式}

\label{idx:115}
你可以把事件看作一种通知，其中包含供被通知的服务处理的额外数据。不过，还有另一种理解方式：事件是状态的变化。想一想，如果你能知道出 bug 时应用逻辑处于什么状态、以及它当时被请求做了什么改动，那么调试应用逻辑中的问题会变得多么容易。这就是事件溯源的一个好处。本质上，它按事件发生的顺序记录所有事件，从而捕获系统中发生的全部变化。

\label{idx:116}
你常常会发现，服务不再需要把自身状态持久化到数据库中，因为把事件存储在系统中的其他地方就足够了。即便确实需要持久化，也可以异步进行。事件溯源带来的另一个好处是完整的审计日志。\emph{图 2.7} 展示了这种架构风格的示意图：

\myGraphic{0.582}{images/B18456_02_07.png}{图 2.7：事件溯源架构}

事件溯源需要投入大量精力来对事件建模，因为每个事件都必须包含还原系统在某一时刻完整状态所需的全部信息。因此，事件一旦存储就不可更改，否则历史记录和状态可能会被破坏或扭曲。这让人联想到 Git 这样的版本控制系统。

\label{idx:117}
定期快照可以缩短发生故障时重放所有事件所需的时间。由于通过重放所有事件来恢复当前状态可能耗时很长，快照允许系统从某个已保存的时间点恢复，并重放该时间点之后发生的事件。快照还有助于适应事件结构随时间发生的变化。

在把事件添加到事件存储、发布事件以及由消费者处理事件之间，存在一定的延迟，而在这期间还可能有更多事件到达。因此，系统状态不会立即更新，但最终会达到一致。

这种设计模式适用于以下情形：需要避免消息代理与数据库之间的耦合；或者不支持涉及多个资源的分布式事务，但又必须让数据库中的原子变更与向消息代理的发布保持一致、同步，以防止出错。

\label{idx:118}
\label{idx:119}
\label{idx:120}
由于对数据同步的需求减少，事件溯源系统结合\textbf{命令查询职责分离}（\textbf{CQRS}）（将在 \emph{第 4 章}《\emph{架构与系统设计模式}》中介绍）往往能提供低延迟，这意味着它们可以几乎无延迟地响应查询。这使得它们很适合\textbf{高频交易}（\textbf{HFT}）系统和活动追踪器等场景。与事件溯源系统不同，\textbf{创建、读取、更新、删除}（\textbf{CRUD}）系统通常采用分层架构，是同步的，并会在更新期间阻塞数据操作，实现起来要简单得多。

没有必要把事件溯源之类的复杂设计模式应用到简单的应用中。这种设计模式适用于以下情形：需要减少或彻底消除数据更新操作之间的冲突；记录已经发生的事件并重放它们，以还原系统的某个状态；回滚更改并保存历史；或者捕捉创建/更改数据的意图、目的或原因。它并不适用于在系统范围内集成数据。

\section{理解分层架构}

\label{idx:121}
如果你的架构开始变得像一团意大利面，或者你想避免这种情况，把组件按层次组织起来可能会有所帮助。还记得模型–视图–控制器（model–view–controller）吗？又或者类似的模式，比如模型–视图–视图模型（model–view–viewmodel）或实体–控制–边界（entity–control–boundary）？这些都是分层架构的典型例子（如果各层在物理上相互分离，也称为 \emph{N} 层架构）。你可以按层次组织代码、创建微服务的层次，或者把这个模式应用到其他你认为能发挥其优势的领域。分层提供了抽象和关注点分离，从而简化开发、测试和依赖注入。这正是引入分层的主要原因。此外，它还有助于在降低复杂性的同时，提升解决方案的模块化程度、可复用性和可维护性。

分层架构往往很容易实现，因为大多数开发者已经了解层次的概念——他们只需开发若干层并把它们堆叠起来，如下图所示。请求和响应逐层传递，直接穿过某一层到达下一层，而每层中的组件只包含该层的业务逻辑：

\myGraphic{0.854}{images/B18456_02_08.png}{图 2.8：一个三层架构的示例，其表示层使用文本界面（图片由 Bartledan 提供，公有领域，来自 Wikimedia Commons：\texthref{https://commons.wikimedia.org/wiki/File:Overview\_of\_a\_three-tier\_application\_vectorVersion.svg}）}

\label{idx:122}
构建高效分层架构的挑战在于，如何明确规定各层之间稳定、定义良好的接口。通常，一个层之上可以叠加多个层。例如，如果你有一个领域逻辑层，它既可以作为表示层的基础，也可以作为向其他服务提供 API 的层的基础。大多数情况下，CRUD 应用都基于这种架构。

分层架构并不局限于单个应用；同样的设计原则也可以扩展到系统层面——例如在微服务中。实现这一点的一种方式是通过服务于前端的后端模式。

\subsection{服务于前端的后端}

\label{idx:123}
多个前端依赖同一个后端，这种情况并不少见。假设你有一个移动应用和一个 Web 应用，两者使用同一个后端。起初这可能是个不错的设计选择。然而，一旦这两个应用的需求和使用场景出现分歧，你的后端就会越来越频繁地需要修改，而且只是为了服务于其中一个前端。这可能导致后端不得不支持相互冲突的需求，比如两种不同的数据存储更新方式，或不同的数据提供场景。与此同时，前端为了正常与后端通信开始需要更多带宽，这也会导致移动应用更耗电。到了这个地步，你就应当考虑为每个前端引入一个独立的后端，只处理该界面特有的需求。

这样一来，你可以把面向用户的应用看作一个包含两层的单一实体：前端和后端。后端可以依赖另一个由下游服务组成的层。请参考下图：

\myGraphic{0.791}{images/B18456_02_09.png}{图 2.9：服务于前端的后端模式}

\label{idx:124}
\label{idx:125}
使用\textbf{服务于前端的后端}（\textbf{BFF}）的缺点在于，为了支持不同的前端客户端，必须重复一些代码。只要这种做法能加快开发，并且长期来看不构成负担，那就可行。但这也意味着，你应当留意把重复的逻辑汇聚到下游服务或可复用库中的机会。有时，仅仅为了汇聚相似的调用而引入一个服务，就能帮助解决重复问题。通常，如果你有很多前端，其中一些仍然可以共用同一个后端，而不会让它产生相互冲突的需求。例如，如果你在为 iOS 和 Android 开发移动应用，可以考虑让它们复用同一个后端，而为 Web 和/或桌面应用使用各自独立的后端。

\section{了解基于模块的架构}

\label{idx:126}
在本节中，\emph{模块}指的是可以在运行时加载和卸载的软件组件，从而提升系统的灵活性。这与 C++20 模块不同，后者旨在解决宏冲突和重复包含问题、提升编译速度并增强封装。关于 C++20 模块，请参见 \emph{第 9 章}《\emph{C++ 代码复用的未来：使用模块}》。

如果你曾经需要让某个组件尽可能少地停机运行，却由于某种原因无法采用常见的容错模式（例如为服务保留冗余副本），那么把这个组件做成基于模块的形式或许能救急。又或者，你只是被这样一种愿景所吸引：一个模块化系统，其中所有模块都有版本管理，并能方便地查找所有可用服务。

\label{idx:127}
其他优势还包括解耦、可测试性和团队协作。正是这些好处，促使人们在 Java 中创建了 OSGi 模块，启发了 SDK macchina.io EDGE，并让它在一些框架中被移植到 C++：Apache Celix 和 C++ Micro Services。采用模块的架构示例包括 Eclipse 这类 IDE、OpenDaylight 这类\textbf{软件定义网络}（\textbf{SDN}）项目，以及 openHAB 这类家庭自动化软件。

\myGraphic{0.539}{images/B18456_02_10.png}{图 2.10：OSGi 架构}

然而，运行时灵活性也是有代价的：大多数框架无法与 OSGi 一起使用。而且，由于启动顺序是动态的，OSGi bundle 的初始化并不简单，框架也可能检测不到不同版本的服务。此外，类加载器需要事先了解框架在运行时所需的资源。因此，大多数项目选择其他方案也就情有可原了。

\label{idx:128}
OSGi 还支持模块之间的自动依赖管理，控制模块的初始化和卸载，以及控制模块的发现。由于它是面向服务的，你可以把使用 OSGi 服务看作类似于在一个“容器”里拥有许多微小的服务。这也是为什么 C++ 的实现之一被命名为 C++ Micro Services。想看看它们的实际效果，请参阅本章末尾链接的 \emph{Getting Started} 指南。

\label{idx:129}
C++ Micro Services 框架以其基于模块的架构而著称，它采用了一个有趣的概念：一种处理单例的新方式。\texttt{Get\allowbreak{}Instance(\allowbreak{})\allowbreak{}} 静态函数不再是简单地返回一个静态实例对象，而是返回一个从 bundle 上下文中获取的服务引用。其核心在于，这个库会动态注册从称为 bundle 的共享库或静态库中加载的服务。因此，单例对象实际上是由非单例类来模拟的，并被你可以配置的服务所取代。这还能让你免于静态反初始化灾难（static deinitialization fiasco）——即多个相互依赖的单例必须以特定顺序卸载。

\label{idx:130}
这个问题也可以用 Alexandrescu 的 Phoenix Singleton 设计模式来解决。它通过使用 placement new 运算符在同一地址上重建已被销毁的对象，来处理死引用问题（即在对象被销毁后仍然访问它）。

\section{介绍六边形架构}

\label{idx:131}
\label{idx:132}
\label{idx:133}
六边形架构也称为端口与适配器。端口定义了适配器与应用核心交互所用的接口，任何适配器都可以在不影响业务逻辑的情况下被替换。这样一来，就能轻松更换外部技术而不影响核心逻辑，甚至可以延后技术决策。

这类架构有助于理解在设计应用时接口是如何以及为何被使用的。六边形架构的核心思想是以这样一种方式来组织应用：使其能够在不依赖外部工具和技术的情况下独立开发和测试，从而把应用的核心逻辑与其所使用的服务分离开来。这种架构在分层架构的基础上引入了端口与适配器的概念。

\begin{myQuote}
使应用能够同样地由用户、程序、自动化测试或批处理脚本驱动，并能够独立于其最终运行时的设备和数据库进行开发和测试。

——Alistair Cockburn，《六边形架构》，2005

\end{myQuote}

\label{idx:134}
\emph{图 2.11} 说明了这种架构风格为何被称为六边形架构。在这里，外部服务通过驱动（主）端口告诉应用该做什么，而应用则使用被驱动（从）端口来完成自己的工作。

\label{idx:135}
驱动（主）适配器把外部输入转换为应用核心逻辑可以理解的格式，被驱动（从）适配器把内部输出转换为外部服务可以理解的格式。因此，适配器把应用与外部世界连接起来。

\myGraphic{0.920}{images/B18456_02_11.png}{图 2.11：六边形架构}

依赖倒置原则（通常用依赖注入模式实现）与适配器模式使驱动适配器和被驱动适配器独立于业务逻辑。\emph{图 2.12} 展示了该原则在六边形架构中是如何应用的。

\myGraphic{0.920}{images/B18456_02_12.png}{图 2.12：六边形架构中的依赖倒置原则与适配器模式}

\label{idx:136}
然而，六边形架构也有其缺点。目前不清楚是仅由主适配器校验输入，还是业务逻辑也参与其中。理想情况下，这些实体应当是相互独立的。此外，复杂的数据库请求可能导致业务逻辑泄漏到适配器中，而事务可能需要复杂的回滚。

这种架构常用于银行类应用，以把处理交易、核对余额和计算手续费等逻辑与用户界面和数据库分离开来。

\section{小结}

在本章中，我们讨论了你能遇到并应用到软件中的各种架构风格。我们讨论了有状态与无状态的服务以及单体架构。我们回顾了 SOA，接着转向微服务，并讨论了扩展应用的方式，特别是使用扩展立方体（scale cube）模型。随后，我们讨论了其他各种架构方法，包括事件驱动架构和基于运行时模块的架构，并展示了在哪里可以发现分层以及分层为何有用。现在你还了解了如何实现事件溯源，并能判断何时使用 BFF 和六边形架构。此外，你现在知道了架构风格如何帮助你实现若干质量属性，以及它们可能带来哪些挑战。

在下一章中，你将学到如何判断哪些属性在给定的系统中是重要的。

\section{问题}

\begin{enumerate}
\item RESTful 服务有哪些特征？
\item 你可以使用哪个工具包来帮助你创建具有韧性的分布式架构？
\item 微服务应该使用集中式存储吗？为什么应该或不应该？
\item 什么时候应该编写有状态的服务而不是无状态的服务？
\item 代理（broker）与中介者（mediator）有何不同？
\item \emph{N} 层（tier）架构与 \emph{N} 层（layer）架构之间有什么区别？
\item 你应该如何着手用基于微服务的架构替换单体应用？
\end{enumerate}

\section{延伸阅读}

\begin{itemize}
\item Robin Flygare 与 Anthon Holmqvist，\emph{Performance Characteristics between Monolithic and Microservice-based Systems}（学位论文）：\href{https://bth.diva-portal.org/smash/record.jsf?pid=diva2\%3A1119785&dswid=-7832}{http://urn.kb.se/resolve?urn=urn:nbn:se:bth-14888}
\item SOA 宣言：\href{https://soa-manifesto.org/default.html}{https://soa-manifesto.org/default.html}
\item Martin Fowler，\emph{Microservices—A Definition of This New Architectural Term}：\href{https://martinfowler.com/articles/microservices.html\#MicroservicesAndSoa }{https://martinfowler.com/articles/microservices.html\#MicroservicesAndSoa}
\item Lewis Coosner，\emph{Microservices and the Scale Cube}，Incus Data Blog：\href{https://incusdata.com/blog/microservices-and-the-scale-cube}{https://incusdata.com/blog/microservices-and-the-scale-cube}
\item C++ Micro Services，\emph{Getting Started}（C++ Micro Services 文档）：\href{http://docs.cppmicroservices.org/en/stable/doc/src/getting_started.html}{http://docs.cppmicroservices.org/en/stable/doc/src/getting\_started.html}
\item Gaurav Gandhi，\emph{Cookie vs Token Authentication}，Naukri Code 360：\href{https://www.naukri.com/code360/library/cookie-vs-token-authentication}{https://www.naukri.com/code360/library/cookie-vs-token-authentication}
\item Teniola Fatunmbi，\emph{A Comparison of Cookies and Tokens for Secure Authentication}，Okta：\href{https://developer.okta.com/blog/2022/02/08/cookies-vs-tokens}{https://developer.okta.com/blog/2022/02/08/cookies-vs-tokens}
\item Robert Engelen，\emph{A Framework for Service-Oriented Computing with C and C++ Web Service Components}，ACM Transactions on Internet Technology，第 8 卷，第 3 期：\href{https://dl.acm.org/doi/10.1145/1361186.1361188}{https://doi.org/10.1145/1361186.1361188}
\item Rahul Awati 与 James Denman，\emph{Sharding}，TechTarget：\href{https://www.techtarget.com/searchoracle/definition/sharding}{https://www.techtarget.com/searchoracle/definition/sharding}
\item PlanetScale，\emph{Sharding vs. Partitioning: What’s the Difference?}：\href{https://planetscale.com/blog/sharding-vs-partitioning-whats-the-difference}{https://planetscale.com/blog/sharding-vs-partitioning-whats-the-difference}
\item Artem Polishchuk，\emph{Scale Cube: A Guide to Scalability in System Architecture}，Bool.dev：\href{https://bool.dev/blog/detail/scale-cube-a-guide-to-scalability-in-system-architecture}{https://bool.dev/blog/detail/scale-cube-a-guide-to-scalability-in-system-architecture}
\item Chris Richardson，\emph{The Scale Cube}，microservices.io：\href{https://microservices.io/articles/scalecube.html}{https://microservices.io/articles/scalecube.html}
\item Robert C. Martin，\emph{The Single Responsibility Principle}，The Clean Code Blog：\href{https://blog.cleancoder.com/uncle-bob/2014/05/08/SingleReponsibilityPrinciple.html}{https://blog.cleancoder.com/uncle-bob/2014/05/08/SingleReponsibilityPrinciple.html}
\end{itemize}

\begin{center}
\begin{tabularx}{\textwidth}{XX}
\toprule
\textbf{立即解锁本书的专属权益} 扫描此二维码，或访问 \href{http://packtpub.com/unlock}{packtpub.com/unlock}，然后按书名搜索本书。 & \begin{center}\includegraphics[width=0.250\linewidth]{images/Unlock_Now.png}\end{center} \\
\midrule
\emph{注意：开始之前请准备好你的购买发票。} &  \\
\bottomrule
\end{tabularx}
\end{center}
